\chapter{Architettura della Data Pipeline (Medallion Architecture)}

\section{Panoramica dell'Architettura}
Per la gestione e la trasformazione dei dati biometrici è stata adottata la **Medallion Architecture**, un pattern architetturale consolidato nel Data Engineering che suddivide la pipeline in tre livelli di qualità progressiva: \textit{Bronze}, \textit{Silver} e \textit{Gold}.

\section{Layer Bronze (Dati Grezzi)}
Il layer **Bronze** rappresenta il punto di ingresso dei dati grezzi estratti dal dataset VitalDB tramite la libreria Python ufficiale.
\begin{itemize}
    \item I tracciati biometrici dei casi selezionati e le tabelle dei dati clinici/di laboratorio vengono scaricati ed estratti nel loro formato originale.
    \item I dati vengono salvati in locale in formato binario compresso Parquet all'interno della directory \texttt{data/bronze/}.
    \item In questa fase non viene applicata alcuna trasformazione o filtraggio, preservando la completa storicizzazione del dato originario.
\end{itemize}

\section{Layer Silver (Data Cleaning e Partizionamento)}
Il layer **Silver** ha il compito di bonificare ed organizzare i dati per renderli idonei alle analisi avanzate.

\subsection{Pulizia dei Dati e Outlier Fisiologici}
I segnali sensoristici grezzi possono presentare artefatti dovuti a perdita di contatto del sensore o disturbi di misurazione. Vengono applicate le seguenti regole di validazione e pulizia:
\begin{itemize}
    \item Rimozione di outlier fisiologicamente non plausibili (es. $\text{HR} < 20$ o $> 250$ bpm, $\text{SpO2} < 50\%$).
    \item Gestione e interpolazione dei valori mancanti (\textit{missing values}) e allineamento dei timestamp.
\end{itemize}

\subsection{Confronto Visivo Prima e Dopo la Pulizia (Bronze vs Silver)}
La Figura \ref{fig:before_after} evidenzia l'impatto concreto del processo di data cleaning su un tracciato reale del dataset:
\begin{itemize}
    \item \textbf{Layer Bronze (Sopra)}: Mostra il segnale frammentato con discontinuità temporali e punti mancanti dovuti a perdite temporanee di contatto dei sensori.
    \item \textbf{Layer Silver (Sotto)}: Mostra lo stesso segnale bonificato tramite interpolazione ed allineamento temporale, risultando in un segnale continuo e coerente per i modelli di Machine Learning.
\end{itemize}

\begin{figure}[h!]
    \centering
    \includegraphics[width=0.88\textwidth]{before_after_cleaning.png}
    \caption{Confronto visivo tra il tracciato della Frequenza Cardiaca grezzo (Bronze) e bonificato (Silver).}
    \label{fig:before_after}
\end{figure}

\subsection{Partizionamento per Reparto Chirurgico}
Per ottimizzare la lettura successiva ed emulare il \textit{partition pruning}, i dati del layer Silver vengono salvati in formato Parquet partizionato fisicamente in sottocartelle basate sul reparto chirurgico di appartenenza del paziente:
$$\texttt{data/silver/department=<REPARTO>/case\_<ID>.parquet}$$

\section{Evoluzione del Formato del Dato: Da Tabella Piatta a BSON Gold}
Oltre alla qualità dei segnali, l'architettura Medallion trasforma la struttura logica del dato. Mentre nel layer Bronze/Silver i dati risiedono in file tabulari piatti, nel layer Gold vengono strutturati in \textbf{documenti BSON gerarchici} contenenti sia il payload sensoristico che la governance del contesto clinico:

\begin{lstlisting}[caption={Trasformazione del dato dal formato Parquet piatto al documento BSON Gold Time Series in MongoDB.}]
// Documento caricato nella collezione 'vital_signals' (Layer Gold)
{
  "_id": ObjectId("65d3f1a0e4b01a2b3c4d5e6f"),
  "timestamp": ISODate("2020-01-01T08:15:30Z"),
  "metadata": {
    "case_id": 1,
    "sensor_name": "Solar8000",
    "department": "General Surgery",
    "age": 64,
    "sex": "M"
  },
  "metrics": {
    "Solar8000_HR": 72.0,
    "Solar8000_PLETH_SPO2": 98.0,
    "Solar8000_NIBP_SBP": 120.0,
    "Solar8000_NIBP_DBP": 80.0,
    "Solar8000_NIBP_MBP": 93.3
  }
}
\end{lstlisting}
